%-------------------------
% Cover Letter in LaTeX
% Classic business-letter layout
%------------------------

\documentclass[letterpaper,11pt]{article}

\usepackage{latexsym}
\usepackage[empty]{fullpage}
\usepackage[hidelinks]{hyperref}
\usepackage[english]{babel}
\input{glyphtounicode}

% Adjust margins
\addtolength{\oddsidemargin}{-0.5in}
\addtolength{\evensidemargin}{-0.5in}
\addtolength{\textwidth}{1in}
\addtolength{\topmargin}{-.75in}
\addtolength{\textheight}{1.5in}

\urlstyle{same}
\raggedbottom
\raggedright
\setlength{\parindent}{0pt}
\setlength{\parskip}{11pt}

\pdfgentounicode=1

\begin{document}
\small

%----------SENDER-----------
Leonard Eckert \\
Nekv\"agen 35 lgh 1011 \\
163 57 Sp\r{a}nga, Sweden

\vspace{18pt}

%----------RECIPIENT-----------
Accenture \\
R\r{a}dmansgatan 42 \\
113 60 Stockholm, Sweden

\vspace{18pt}

\hfill\today

\vspace{18pt}

\textbf{Application: Spring Internship 2027 -- AI \& Data Track}

\vspace{6pt}

Dear Hiring Team,

I am applying for the AI \& Data track of Accenture's Spring Internship. I started my Master's in AI and Language because I think this is a genuinely exciting time to be studying the field: language models have gone from research demos to things companies actually build products around in just a couple of years, and I want to understand that shift well enough to help build with it, not just use it. What draws me to Accenture specifically is the consulting model itself: sitting with clients and translating that technology into decisions they can actually make, not leaving it as a research result.

Over eight internships at IBM, I worked in close contact with clients. In Client Engineering I deployed a demo cluster on OpenShift during a client engagement in Stuttgart, prepared workshop materials in Singapore, and built an agent-based demo with watsonx Orchestrate for partner enablement in Vienna. At IBM Consulting I joined an SAP project for the German Navy on material readiness, where I saw how much careful stakeholder alignment a large organisation needs before technology changes anything. Across all of these, I learned that a good recommendation depends as much on explaining data clearly to people who are not data people as on the analysis itself.

On the technical side, I currently work as a Data Engineer at Bosch Research, where I built an Apache Spark application on Databricks that processes over 30 billion records of vehicle sensor data, and I support platform work in Azure. Before that, at IBM Research in Dublin, I designed an evaluation framework with tailored metrics for large language models, which tackles a question every AI client eventually asks: is this model actually good enough for our use case? My Bachelor's thesis on a semantic alignment agent for natural language generation came out of the same work.

Many of Accenture's AI and data projects are built on Databricks and Azure, the same platforms I work with every day, so I could contribute from the first week. I am now based in Stockholm for my Master's at Stockholm University and speak Swedish at an advanced level alongside German and English, which would let me work with Swedish clients as well as international teams.

I would welcome the opportunity to discuss my application further.

\vspace{18pt}

Best regards, \\
Leonard Eckert

\end{document}
